\section{Approach Details}
\label{SEC:approach}

\begin{figure}[t]
  \center{}
  \fbox{\includegraphics[scale=.5]{approach}}
  \caption{Diagram illustrating CrashSimulator's approach}
  \label{figure:approach}
\end{figure}

In order to understand CrashSimulator's operation it is necessary to have a
more detailed understanding of how its components fit together and what
steps are required to use it in testing an application.
Figure~\ref{figure:approach} illustrates the steps required to gather
the anomaly information CrashSimulator requires, how that information is to
be transformed into mutated traces and what output CrashSimulator yields
once its replay and analysis process completes.

\subsection{Anomaly Identification} \label{subsec:anomalyidentification}

The first step in CrashSimulator's testing process is identifying an
environmental anomaly to test an application against.  This can be
accomplished either manually or by examining results of other applications
running in the target environment with existing tools.  For example, this
work utilized the results of both NetCheck~\cite{Zhuang_NSDI_2014} and
CheckAPI~\cite{rasley2015detecting}, as well as manual analysis, to develop
the set of anomalies used in our evaluation.  NetCheck, which can be
thought of as a precursor to CrashSimulator, is a tool that can determine
the cause of a failure in a networked application, without any specific
knowledge of the application or its language. It is readily able to
identify a wide variety of anomalous behaviors related to network
communication between one or more hosts.  CheckAPI performs a similar
function by comparing the behavior of an API implementation to a reference
model.

Another source of anomalies is the public bug trackers of projects
available on the web.  If the presence of a bug is indicated by the results
and side effects of a series of system calls then CrashSimulator can be
used to inject these same results and side effects at appropriate times in
other applications.  This is ideal when determining whether or not an
application is vulnerable to a widely publicized bug.  Similarly, bugs can
be taken from similar projects in order to expand a project's set of
available anomalies.

In selecting anomalies for use it is important to understand that they can
be as subtle as a single configuration flag being set differently under
specific circumstances.  For example, in BSD and OS X sockets returned by
the {\tt accept()} system call inherit the values of the {\tt O\_NONBLOCK}
and {\tt O\_ASYNC} flags from their parent socket.  In Linux, these values
are not inherited.  As a result, the correct behavior is to always
explicitly set these flags to the appropriate values on child sockets.
This environmental difference presented itself as a bug in the socket
implementation of Python 3.2.  Python took advantage of non-blocking
sockets in order to implement timeout capabilities for sockets.  Because
Python did not specifically set the {\tt O\_NONBLOCK} flag on the child
sockets underlying its socket object abstraction, it was possible that,
depending on the environment, the Python socket object would report the
socket as either blocking or non-blocking while the actual underlying
socket was configured for the opposite behavior.

Anomalies can also be found around more complex operations.  Moving files
around in a filesystem is an extremely common operation that, on face
value, seems to fairly simplistic to carry out.  In reality, things become
complex quickly once a few implementation details are taken into account:
namely limitations in Linux's {\tt rename()} system call and the way
Linux's directory hierarchy can abstract away complicated device
configurations.

\preston{This example might need to be changed.  One of the pieces of
feedback we received from reviewers was that some of our bugs could be
detected without replaying and injecting.  That is, they could be
identified by analyzing a dead strace recording from an application.  This
example will be much stronger if we inject rename() = -1 EXDEV and then
look for this pattern in the application's response once we have that
machinery in place.}

Trivial cases of moving a file (i.e. those where the source and destination
are on the same device) are accomplished using the {\tt rename()} system
call.  In cases where the source and destination are not on the same device
this system call fails and the application must handle the operation
manually. This means that \emph{any} application that moves files around on
a Linux system must be able to deal with the situation where some part of
the directory structure may physically reside on a different device.
During our evaluation, we modeled what it takes to move a file correctly
using the GNU Coreutils {\tt mv} command as our ``gold standard.''  We
found many applications do not handle the above situation correctly
including the {\tt shutils} module from the Python 2.7.9 standard library.
It did not correctly preserve timestamps across the move, it did not
correctly apply extended file attributes to the destination file after the
move, and it did not verify the file's inode number between initially
examining the source file and when the copy was initiated (a common cause
of race conditions).  Failing to perform these operations exposes
applications relying on {\tt shutils} to misordered processing, data loss,
or security issues.

\subsection{Converting Anomalies to System Call Sequences}

%  This paragraph should cover the "describe manual effort" feedback
Once a new anomaly has been identified, it is distilled into a sequence of
system calls that respond to the anomaly.  This process involves examining
the system call traces in order to identify a series of calls indicative of
an inappropriate response to a specified anomaly.  This analysis should
identify any mutations that may need to be made to the trace
to simulate the anomalous environment (i.e. what modifications must be made
to system call return values and program state during replay).  Much like
the effort involved in constructing a unit test to check for correct
behavior in a piece of code, this process involves an initial outlay of
user skill and effort that will be paid back as it is used repeatedly over
time to test many different applications.

As a more concrete example of the above, consider an anomaly
involving unusual file types we will later address in our evaluation of
CrashSimulator.  This anomaly appears when a call to {\tt stat()} or a similar system
call returns a structure with a {\tt st\_mode} member containing an unexpected
value. Consider the line of {\tt strace} output representing a call to {\tt
  fstat()}:
\begin{quote}
  {\tt 8936  fstat64(3, \{st\_dev=makedev(0, 40), st\_ino=54993216, st\_mode=S\_IFREG ...\}) = 0}
\end{quote}
The third member of the returned structure indicates that the file is a
regular file by showing that {\tt st\_mode} flag is set in the {\tt st\_mode}
member.  CrashSimulator can mutate this  line to the following:

\begin{quote}
  {\tt 8936  fstat64(3, \{st\_dev=makedev(0, 40), st\_ino=54993216, st\_mode=S\_IFCHR ...\}) = 0}
\end{quote}

The trace containing this modified line can then be replayed in order to
test how the application responds when a file that is expected to be
regular is actually a character device. In our evaluation we further
explore the behavior of applications as they encounter other possible valid
values of {\tt st\_mode} in the course of testing.


\subsection{Trace Recording and Anomaly Injection}

\preston{This is worded strangely}
The trace into which anomalies will be injected is generated by recording
the application being tested as it executes in some environment where the
chosen anomaly isn't present.

Once this trace has been generated, CrashSimulator's replay-based approach
makes it easy to inject anomalies into an execution by introducing the
features of the anomaly into the trace.  Because replaying the modified
trace faithfully follows the described system call behavior, the
application is exposed to the anomalous conditions that were injected.  For
example, suppose an anomaly causes a call to {\tt read()} to fail with
return value of -1 and {\tt errno} set to {\tt EINVAL}.  Exposing an
application to this situation only requires editing the system call trace,
such that the line describing this read call reflects the above values.

Our implementation of CrashSimulator uses a combination of {\tt rr}, a
record-and-replay debugger, and our own CrashSimulator process supervisor
to accomplish the above.  We rely on {\tt rr's} trace recording to format
to store the information necessary to perform the bulk of uninteresting
replay and {\tt strace}-style output to describe user facing system call
sequences that will be mutated during testing.  {\tt Strace} is a good fit
for this purpose because it provides a familiar way to describe an
application's system call activity in terms of a user-friendly
representation of the inputs, outputs, and side effects of each system call
an application makes.  This means CrashSimulator's user can easily modify
an {\tt strace} recording using a text editor in order to express an
anomaly in terms of system call results and side effects.  These
modifications can typically be automated with the help of a "mutator"
script (and were for the anomalies used in our evaluation).

\subsection{Analyzing an Execution}

CrashSimulator's determination about whether an application may be
responding to an anomaly comes down to whether the application alters its
system call behavior or not.  When CrashSimulator injects anomalous
behavior into a replay execution it is expected that the application under
test will deal with this anomaly in some way. For example, applications can
be expected to behave differently when the structure returned by a call to
{\tt stat()} indicates the target file is a symlink, rather than a regular
file.  CrashSimulator bases its assessment on the assumption that an
application's efforts to deal with the anomaly will result in it executing
different program paths (and therefore different system calls) than were
executed when the system call trace was recorded.  If the application being
replayed does not alter its behavior to deal with the anomaly, it has not
correctly handled this flaw --- a failing result.  Alternatively, if the
application does deviate, it is likely that the application is taking some
action to handle the injected condition --- an indication of possibly
correct behavior.

This approach is often sufficient to classify application behavior.  As
with traditional tests, CrashSimulator does not tell us that an application
is bug free in a given situation; instead, it asserts that an application
has incorrectly handled a given situation.  In our experience, actually
exposing applications to anomalies CrashSimulator reported as unhandled
resulted in bugs ranging from simple hangs to more dramatic crashes that
had the side effect of consuming all available resources or filling
available disk space with garbage.
